% Part:sets-functions-relations
% Chapter: size-of-sets
% Section: enumerations

\documentclass[../../../include/open-logic-section]{subfiles}

\begin{document}

\olfileid{sfr}{siz}{enm}

\olsection{లెక్కింపులు, \usetoken{S}{enumerable} సమితులు}

\begin{editorial}
  ఈ విభాగం సమితుల లెక్కింపులను చర్చిస్తూ, వాటిని $\PosInt$ నుంచి
  సంగ్రస్త ప్రమేయాలుగా నిర్వచిస్తుంది. గణిత నేపథ్యం తక్కువగా ఉన్న
  పాఠకుల కోసం విషయాన్ని నెమ్మదిగా వివరిస్తుంది. లెక్కింపులను వేరుగా,
  అంటే $\Nat$తో (లేదా దాని ఆరంభ ఖండంతో) ద్విగుణ ప్రమేయాలుగా
  నిర్వచించే సంక్షిప్త ప్రత్యామ్నాయ రూపం \olref[enm-alt]{sec}లో ఉంది.
\end{editorial}

\begin{explain}
సమితిలోని !!{element}sను జాబితాగా రాసి సమితులకు ఇప్పటికే ఉదాహరణలు
ఇచ్చాం. ఒక సమితి అనంతమైనా, దానిలోని !!{element}sను ఎలా, ఎప్పుడు
జాబితా చేయగలమో ఇప్పుడు మరింత సాధారణంగా చర్చిద్దాం.
\end{explain}

\begin{defn}[లెక్కింపు: అనౌపచారికంగా]
అనౌపచారికంగా, $A$ సమితి యొక్క \emph{లెక్కింపు} అంటే $A$లోని
!!{element}sతో కూడిన (అనంతం కావచ్చుననే) జాబితా; $A$లోని ప్రతి
!!{element} ఆ జాబితాలో ఏదో ఒక పరిమిత స్థానంలో కనిపించాలి.
$A$కు ఒక లెక్కింపు ఉంటే, $A$ను \emph{!!{enumerable}} అంటాం.
\end{defn}

\begin{explain}
లెక్కింపుల గురించి కొన్ని విషయాలు:
\begin{enumerate}
\item మొదలు ఉండి, మొదటిది కాని ప్రతి !!{element}కు తనకు వెంటనే ముందు
  ఒకే ఒక్క !!{element} ఉండే జాబితాలనే లెక్కింపులుగా పరిగణిస్తాం.
  మరోలా చెప్పాలంటే, జాబితాలోని మొదటి !!{element}కు, మరే ఇతర
  !!{element}కు మధ్య పరిమిత సంఖ్యలోనే మూలకాలు ఉంటాయి. ప్రత్యేకించి,
  లెక్కింపులోని ప్రతి !!{element}కు పరిమిత స్థానం ఉంటుంది: మొదటి
  !!{element} స్థానం $1$, రెండవదాని స్థానం $2$, ఇలాగే కొనసాగుతుంది.
\item ఒకే సమితి $A$కు !!{element}s కనిపించే క్రమం వేరుగా ఉండే
  లెక్కింపులు ఉండవచ్చు: $4$, $1$, $25$, $16$,~$9$ అనే జాబితా మొదటి
  ఐదు వర్గ సంఖ్యల (సమితిని) $1$, $4$, $9$, $16$,~$25$ ఎంత బాగా
  లెక్కిస్తుందో అంతే బాగా లెక్కిస్తుంది.
\item పునరుక్తులు గల లెక్కింపులూ లెక్కింపులే: $1$, $1$, $2$, $2$,
  $3$, $3$,~\dots{} అనేది $1$, $2$, $3$,~\dots{} లెక్కించే సమితినే
  లెక్కిస్తుంది.
\item ఒక లెక్కింపును \emph{పేర్కొన్నప్పుడు} క్రమం, పునరుక్తులు
  ముఖ్యమే: ధన పూర్ణ సంఖ్యల లెక్కింపును $1$, $2$, $3$, $1$, \dots{}
  అని మొదలుపెట్టవచ్చు; కానీ $1$, $2$, $3$, $4$,~\dots అనే ప్రామాణిక
  విధంగా లెక్కిస్తే నమూనా మరింత సులభంగా కనిపిస్తుంది.
\item లెక్కింపులకు తప్పనిసరిగా మొదలు ఉండాలి: \dots, $3$, $2$, $1$
  అనేది ధన పూర్ణ సంఖ్యల లెక్కింపు కాదు, ఎందుకంటే దానికి మొదటి
  !!{element} లేదు. అనౌపచారిక నిర్వచనం నుంచి ఇది ఎలా వస్తుందో
  చూడటానికి, ``జాబితాలో 76 ఏ స్థానంలో కనిపిస్తుంది?'' అని ఆలోచించండి.
\item కిందిది ధన పూర్ణ సంఖ్యల లెక్కింపు కాదు:
  $1$, $3$, $5$, \dots, $2$, $4$, $6$, \dots\@ సమస్య ఏమిటంటే సరి
  సంఖ్యలు పరిమిత స్థానాల్లో కాకుండా $\infty + 1$, $\infty + 2$,
  $\infty + 3$ స్థానాల్లో వస్తాయి.
\item శూన్య సమితి లెక్కించదగినది: దాన్ని ఖాళీ జాబితా లెక్కిస్తుంది!{}
\end{enumerate}
\end{explain}

\begin{prop}
  $A$కు ఒక లెక్కింపు ఉంటే, పునరుక్తులు లేని లెక్కింపు కూడా ఉంటుంది.
\end{prop}

\begin{proof}
  ప్రతి $x_i$, $A$లోని !!{element} అయ్యేలా $A$కు $x_1$, $x_2$,
  \dots{} అనే లెక్కింపు ఉందనుకుందాం. పునరావృతమైన !!{element}sను
  తొలగించి లెక్కింపులోని పునరుక్తులను తీసేయవచ్చు. ఉదాహరణకు, $x_i$
  అనేది $x_1$, \dots, $x_{i-1}$లో లేని $A$ యొక్క !!a{element} అయితే
  దాన్ని కొత్త జాబితాలో ఉంచాలి; ఇప్పటికే $x_1$, \dots,~$x_{i-1}$లో
  ఉంటే $x_i$ను జాబితా నుంచి తొలగించాలి.
\end{proof}

లెక్కింపులనూ !!{enumerable} సమితులనూ బాగా అర్థం చేసుకుని వాటి గురించి
విషయాలను నిరూపించాలంటే మరింత కచ్చితమైన నిర్వచనం అవసరమని చివరి వాదన
చూపుతుంది. కిందిది ఆ నిర్వచనాన్ని ఇస్తుంది.

\begin{defn}[లెక్కింపు: ఔపచారికంగా]
$A \neq \emptyset$ సమితి యొక్క \emph{లెక్కింపు} అంటే ఏదైనా
!!{surjective} ప్రమేయం $f \colon \PosInt \to A$.
\end{defn}

\begin{explain}
ఔపచారిక నిర్వచనం, అనంతం కావచ్చుననే జాబితాను వాడిన అనౌపచారిక
నిర్వచనం సమానమని నిర్ధారించుకుందాం. మొదట, $\PosInt$ నుంచి $A$
సమితికి వెళ్లే ప్రతి !!{surjective} ప్రమేయం $A$ను లెక్కిస్తుంది.
ఇలాంటి ప్రమేయం పైన అనౌపచారికంగా నిర్వచించిన లెక్కింపును నిర్ణయిస్తుంది:
$f(1)$, $f(2)$, $f(3)$, \dots అనే జాబితా. $f$ !!{surjective} కనుక,
$A$లోని ప్రతి !!{element} ఏదో ఒక $n \in \PosInt$కు $f(n)$ విలువ
అవుతుంది. అందువల్ల $A$లోని ప్రతి !!{element} జాబితాలో ఏదో ఒక పరిమిత
స్థానంలో కనిపిస్తుంది. ప్రమేయం !!{injective} కాకపోవచ్చు కనుక జాబితాలో
పునరుక్తులు ఉండవచ్చు; పైన గమనించినట్లు అది అనుమతించబడుతుంది.

మరోవైపు, $A$లోని !!{element}s అన్నింటినీ లెక్కించే జాబితా ఇచ్చినప్పుడు,
$f(n)$ను జాబితాలోని $n$వ !!{element}గా---లేదా $n$వ !!{element} లేకపోతే
జాబితాలోని చివరి !!{element}గా---తీసుకుని !!a{surjective} ప్రమేయం
$f\colon \PosInt \to A$ను నిర్వచించవచ్చు. $A$ ఖాళీగా ఉండి, అందువల్ల
జాబితా ఖాళీగా ఉన్న సందర్భంలో మాత్రమే ఇది !!a{surjective} ప్రమేయాన్ని
ఇవ్వదు. కాబట్టి ప్రతి ఖాళీ కాని జాబితా ఒక !!a{surjective} ప్రమేయం
$f\colon \PosInt \to A$ను నిర్ణయిస్తుంది.
\end{explain}

\begin{defn}
  \ollabel{defn:enumerable}
  $A$ సమితి ఖాళీగా ఉంటే లేదా దానికి లెక్కింపు ఉంటేనే, అప్పుడే మాత్రమే
  అది !!{enumerable}.
\end{defn}

\begin{ex}
ధన పూర్ణ సంఖ్యలను ($\PosInt$) లెక్కించే ప్రమేయం $f(n) = n$గా ఇచ్చిన
తాదాత్మ్య ప్రమేయమే. సహజ సంఖ్యలు $\Nat$ను లెక్కించే ప్రమేయం
$g(n) = n - 1$.
\end{ex}

\begin{ex}
\begin{align*}
f(n) & = 2n \text{ మరియు}\\
g(n) & = 2n - 1
\end{align*}
అని ఇచ్చిన $f\colon \PosInt \to \PosInt$, $g \colon \PosInt \to
\PosInt$ ప్రమేయాలు వరుసగా ధన సరి పూర్ణ సంఖ్యలనూ, ధన బేసి పూర్ణ
సంఖ్యలనూ లెక్కిస్తాయి. అయితే ఏ ప్రమేయమూ !!{surjective} కాదు కనుక,
రెండింటిలో ఏదీ $\PosInt$కు లెక్కింపు కాదు.
\end{ex}

\begin{prob}
ధన వర్గ సంఖ్యలు $1$, $4$, $9$, $16$, \dots కు ఒక లెక్కింపును
నిర్వచించండి.
\end{prob}

\begin{ex}
$\lceil x \rceil$ అనేది $x$ను దానికి సమీపంలోని పై పూర్ణ సంఖ్యకు
చేర్చే \emph{పై పూర్ణాంక} ప్రమేయం అనుకుందాం. అప్పుడు
$f(n) = (-1)^{n} \lceil \frac{(n-1)}{2}\rceil$ ప్రమేయం పూర్ణ సంఖ్యల
సమితి $\Int$ను లెక్కిస్తుంది. $\Int$ విలువలను $f$ ధన, ఋణ పూర్ణ
సంఖ్యల మధ్య ``అటూ ఇటూ దూకుతూ'' ఎలా ఉత్పత్తి చేస్తుందో గమనించండి:
\[
\begin{array}{c c c c c c c c}
f(1) & f(2) & f(3) & f(4) & f(5) & f(6) & f(7) & \dots \\ \\
- \lceil \tfrac{0}{2} \rceil & \lceil \tfrac{1}{2}\rceil & - \lceil \tfrac{2}{2} \rceil & \lceil \tfrac{3}{2} \rceil & - \lceil \tfrac{4}{2} \rceil  & \lceil \tfrac{5}{2}
\rceil & - \lceil \tfrac{6}{2} \rceil & \dots \\ \\
0 & 1 & -1 & 2 & -2 & 3 & -3 & \dots
\end{array}
\]
\sourcecorrection{OLSIZ-001}{ఆంగ్ల మూలపు పట్టికలో $f(7)$ మరియు
$-\lceil 6/2\rceil$ ఉన్నా, విలువల వరుసలో దానికి సరిపడే $-3$ లేదు.
ఆ నియంత్రక సూత్రం ప్రకారం $-3$ను చేర్చాం.}
$f$ను సందర్భాలవారీగా ఇలా నిర్వచించినట్లు కూడా భావించవచ్చు:
\[
f(n) = \begin{cases}
  0 & \text{$n = 1$ అయితే}\\
  n/2 & \text{$n$ సరి సంఖ్య అయితే}\\
  -(n-1)/2 & \text{$n$ బేసి సంఖ్య మరియు $>1$ అయితే}
  \end{cases}
\]
\end{ex}

\begin{prob}
  $A$, $B$లు !!{enumerable} అయితే $A \cup B$ కూడా లెక్కించదగినదని
  చూపండి. ఇందుకోసం !!{surjective} ప్రమేయాలు $f\colon \PosInt \to A$,
  $g\colon \PosInt \to B$ ఉన్నాయని అనుకుని, !!a{surjective} ప్రమేయం
  $h\colon \PosInt \to A \cup B$ను నిర్వచించి అది !!{surjective} అని
  నిరూపించండి. $A$ లేదా~$B = \emptyset$ అయ్యే సందర్భాలనూ పరిశీలించండి.
\end{prob}

\begin{prob}
  $B \subseteq A$, $A$ !!{enumerable} అయితే $B$ కూడా లెక్కించదగినది
  అని చూపండి. ఇందుకోసం !!a{surjective} ప్రమేయం $f\colon \PosInt \to
  A$ ఉందనుకోండి. !!a{surjective} ప్రమేయం $g\colon \PosInt \to B$ను
  నిర్వచించి, అది !!{surjective} అని నిరూపించండి. $B = \emptyset$
  అయితే ఏమవుతుంది?
\end{prob}

\begin{prob}
  $A_1$, $A_2$, \dots, $A_n$ అన్నీ !!{enumerable} అయితే
  $A_1 \cup \dots \cup A_n$ కూడా లెక్కించదగినదని $n$పై ఆగమనం
  ద్వారా చూపండి. రెండు సమితులు $A$,~$B$ !!{enumerable} అయితే
  $A \cup B$ కూడా లెక్కించదగినది అనే ఫలితాన్ని వాడుకోవచ్చు.
\end{prob}

ఒక సమితిలోని !!{element}sను జాబితా చేసేటప్పుడు మొదటి !!{element}ను
$1$వదిగా లెక్కించడం మరింత సహజంగా అనిపించవచ్చు. కానీ గణితశాస్త్రజ్ఞులు
వస్తువులను లెక్కించడానికి సహజ సంఖ్యలు $\Nat$ను వాడటానికి ఇష్టపడతారు.
వారు జాబితాలోని $0$వ, $1$వ, $2$వ మొదలైన !!{element}s గురించి
మాట్లాడతారు. దానికి అనుగుణంగా, $\Nat$ నుంచి $A$కు వెళ్లే
!!a{surjective} ప్రమేయంగా లెక్కింపును నిర్వచించవచ్చు. ఈ రెండు
నిర్వచనాలు సమానమే.

\begin{prop}\ollabel{prop:enum-shift}
  !!a{surjection} $f\colon \PosInt \to A$ ఉంటేనే, అప్పుడే మాత్రమే
  !!a{surjection} $g\colon \Nat \to A$ ఉంటుంది.
\end{prop}

\begin{proof}
  !!a{surjection} $f\colon \PosInt \to A$ ఇచ్చినప్పుడు, ప్రతి
  $n \in \Nat$కు $g(n) = f(n+1)$గా నిర్వచించవచ్చు. $g\colon \Nat
  \to A$ !!{surjective} అని సులభంగా చూడవచ్చు. తిరిగి,
  !!a{surjection} $g\colon \Nat \to A$ ఇచ్చినప్పుడు
  $f(n) = g(n-1)$గా నిర్వచించండి.
\end{proof}

దీంతో మనకు కింది ఫలితం లభిస్తుంది:

\begin{cor}\ollabel{cor:enum-nat}
$A$ సమితి ఖాళీగా ఉంటే లేదా !!a{surjective} ప్రమేయం
$f\colon \Nat \to A$ ఉంటేనే, అప్పుడే మాత్రమే అది !!{enumerable}.
\end{cor}

$A$ సమితిలోని !!{element}s జాబితాను పునరుక్తులు లేని జాబితాగా
మార్చవచ్చని పైన చర్చించాం. లెక్కింపుల విషయంలోనూ ఇది నిజమే; కానీ
కచ్చితంగా సూత్రీకరించి నిరూపించడం కొంచెం కష్టం. ప్రతి $n \in
\PosInt$ వద్ద $f\colon \PosInt \to A$ ప్రమేయం నిర్వచితమై ఉండాలి.
$A$లో పరిమిత సంఖ్యలో మాత్రమే !!{element}s ఉంటే, అనంత సంఖ్యలో
!!{element}s గల $\PosInt$పై నిర్వచితమై, $A$లోని !!{element}s
అన్నిటినీ విలువలుగా తీసుకుంటూ ఏ విలువనూ రెండుసార్లు తీసుకోని ప్రమేయం
ఉండటం అసాధ్యం. ఆ సందర్భంలో, అంటే పునరుక్తులు లేని జాబితా పరిమితంగా
ఉన్నప్పుడు, $f$కు పరిమిత సంఖ్యలోనే !!{element}s గల వేరొక ప్రవేశాన్ని
ఎంచుకోవాలి. పునరుక్తులు లేకపోవడం అంటే $f$ !!{injective} కావాలి.
అది !!{surjective} కూడా కనుక, ఏదో ఒక పరిమిత సమితి
$\{1, \dots, n\}$ లేదా $\PosInt$కు, $A$కు మధ్య !!a{bijection} కోసం
చూస్తున్నాం.

\begin{prop}\ollabel{prop:enum-bij}
$f\colon \PosInt \to A$ !!{surjective} అయితే (అంటే $A$కు లెక్కింపు
అయితే), $Z$ అనేది $\PosInt$ లేదా ఏదో ఒక $n \in \PosInt$కు
$\{1, \dots, n\}$ అయ్యేలా !!a{bijection} $g\colon Z \to A$ ఉంటుంది.
\end{prop}

\begin{proof}
  $g$ ప్రమేయాన్ని పునరావృత్తంగా నిర్వచిస్తాం: $g(1) = f(1)$గా
  తీసుకోండి. $g(i)$ ఇప్పటికే నిర్వచితమైతే, $f(1)$, $f(2)$, \dots{}లో
  $g(1)$, \dots, $g(i)$ మధ్య ఇంకా లేని మొదటి విలువ ఉంటే, దానిని
  $g(i+1)$గా తీసుకోండి. $A$లో కేవలం $n$ !!{element}s ఉంటే,
  $g(1)$, \dots, $g(n)$ అన్నీ నిర్వచితమవుతాయి; అప్పుడు
  $g\colon \{1, \dots, n\} \to A$ ప్రమేయాన్ని నిర్వచించినట్లే.
  $A$లో అనంత సంఖ్యలో !!{element}s ఉంటే, ఏ $i$కైనా $f(1)$, $f(2)$,
  \dots లెక్కింపులో $g(1)$, \dots, $g(i)$ మధ్య ఇంకా లేని
  $A$ యొక్క !!a{element} తప్పక ఉంటుంది. ఈ సందర్భంలో
  $g\colon \PosInt \to A$ ప్రమేయాన్ని నిర్వచించాం.

  $A$ యొక్క ఏ మూలకమైనా $f(1)$, $f(2)$, \dots{}లో ఉంటుంది ($f$
  !!{surjective} కనుక); అందుచేత ఏదో ఒక $i$కు చివరకు $g(i)$ విలువగా
  వస్తుంది. కాబట్టి $g$ !!{surjective}. అది !!{injective} కూడా:
  $j < i$ అయ్యే jకు $g(j) = g(i)$ అయితే, $g(i)$ ఇప్పటికే
  $g(1)$, \dots, $g(i-1)$లో ఉండేది; అది $g$ను నిర్వచించిన విధానానికి
  విరుద్ధం.
\end{proof}

\begin{cor}\ollabel{cor:enum-nat-bij}
$A$ సమితి ఖాళీగా ఉంటే లేదా $N = \Nat$ గానీ, ఏదో ఒక $n \in \Nat$కు
$N = \{0, \dots, n\}$ గానీ అయ్యేలా !!a{bijection}
$f\colon N \to A$ ఉంటేనే, అప్పుడే మాత్రమే అది !!{enumerable}.
\end{cor}

\begin{proof}
$A$ !!{enumerable} అవడం అంటే $A$ ఖాళీగా ఉండటం లేదా !!a{surjective}
$f\colon \PosInt \to A$ ఉండటం. \olref{prop:enum-bij} ప్రకారం, రెండవది
$Z = \PosInt$ లేదా ఏదో ఒక $n \in \PosInt$కు
$Z = \{1, \dots, n\}$ అయ్యేలా !!a{bijective} ప్రమేయం
$f\colon Z \to A$ ఉండటంతో సమానం. \olref{prop:enum-shift} నిరూపణలోని
వాదన ప్రకారమే, అది తిరిగి $N = \Nat$ లేదా
$N = \{0, \dots, n-1\}$ అయ్యేలా !!a{bijection} $g\colon N \to A$
ఉండటంతో సమానం.
\end{proof}

\begin{prob}
  \olref[sfr][siz][enm]{defn:enumerable} ప్రకారం, $A = \emptyset$
  లేదా !!a{surjective} ప్రమేయం $f\colon \PosInt \to A$ ఉంటేనే,
  అప్పుడే మాత్రమే $A$ సమితి లెక్కించదగినది. ``!!{enumerable} సమితి''ని
  కచ్చితంగా మరో విధంగా కూడా నిర్వచించవచ్చు: !!a{injective} ప్రమేయం
  $g\colon A \to \PosInt$ ఉంటేనే, అప్పుడే మాత్రమే సమితి లెక్కించదగినది.
  ఈ నిర్వచనాలు సమానమని చూపండి; అంటే, $A = \emptyset$ కావడం లేదా
  !!a{surjective} ప్రమేయం $f\colon \PosInt \to A$ ఉండటం, !!a{injective}
  ప్రమేయం $g\colon A \to \PosInt$ ఉండటంతో సమానమని చూపండి.
\end{prob}

\end{document}
